% Zone „kennst du schon“ – Fokus Nullstellen (quadratische-funktionen.md, Einheit 4)
\setcounter{aufgabe}{0}
\einheitenkopf{Das kennst du schon}

\begin{aufgabe}{Ich kann Zahlen quadrieren, auch negative Zahlen und Dezimalzahlen. Berechne ohne Taschenrechner.}
\begin{teilezwei}
\tz $(-5)^2 =$ \leerfeld & \tz $(1{,}5)^2 =$ \leerfeld \\
\end{teilezwei}
\end{aufgabe}

\begin{aufgabe}{Ich kann berechnen, wo eine Gerade die x-Achse schneidet. Berechne die Stelle $x_0$, an der die Gerade die x-Achse schneidet.}
\nullstellenliste{2x - 8, -3x - 6}
\end{aufgabe}

\begin{aufgabe}{Ich kann eine Klammer mit der binomischen Formel auflösen. Löse die Klammer auf und fasse zusammen.}
\begin{gleichungsraster}[2]
\gl{(x + 5)^2} & \gl{(x - 3)^2 - 5} \\
\end{gleichungsraster}
\end{aufgabe}

\begin{aufgabe}{Ich kann Quadratwurzeln ziehen und weiß, wann es keine gibt.}
\begin{teile}
\teil Ziehe die Wurzel mit dem Taschenrechner und runde auf zwei Nachkommastellen. \\ $\sqrt{7} \approx$ \leerfeld
\teil Entscheide: Gibt es eine Zahl, die man für $\sqrt{-4}$ schreiben kann? \quad \janein
\end{teile}
\end{aufgabe}

\begin{aufgabe}{Ich kann eine quadratische Gleichung lösen. Löse die Gleichung und gib alle Lösungen an. Bei b) nimmst du die p-q-Formel.}
\begin{gleichungsraster}[3]
\gl{x^2 = 64} & \gl{x^2 + 4x - 12 = 0} \\
\end{gleichungsraster}
\end{aufgabe}
